% 主小行星带（综述）
% license CCBYSA3
% type Wiki

本文根据 CC-BY-SA 协议转载翻译自维基百科\href{https://en.wikipedia.org/wiki/Asteroid_belt}{相关文章}。

\begin{figure}[ht]
\centering
\includegraphics[width=6cm]{./figures/f9c18c4daeef0b22.png}
\caption{} \label{fig_belt_1}
\end{figure}
\begin{figure}[ht]
\centering
\includegraphics[width=6cm]{./figures/cf13395dc35061d2.png}
\caption{到目前为止，腰带内最大的物体是矮行星谷神星。小行星带的总质量明显小于冥王星大约是冥王星月亮的两倍卡隆。} \label{fig_belt_2}
\end{figure}
小行星带是太阳系中的一个圆环形区域，以太阳为中心，大致横跨行星木星和火星轨道之间的空间。该区域包含大量固态、不规则形状的天体，称为小行星或小行星状星体。这些已被确认的天体大小各异，但都远小于行星，且平均彼此相距约一百万千米（约六十万英里）。这个小行星带也被称为主小行星带或主带，以区别于太阳系中其他小行星族群$^{[1]}$。

小行星带是太阳系中最小且最内侧的环星盘。位于其他区域的小型太阳系天体类别，还包括近地天体、半人马天体、柯伊伯带天体、散射盘天体、塞德纳型天体以及奥尔特云天体。主带总质量中约 60\% 集中在四颗最大的小行星：谷神星（Ceres）、灶神星（Vesta）、智神星（Pallas）和健神星（Hygiea）中。小行星带的总质量估计约为月球质量的 3\%$^{[2]}$。

谷神星是小行星带中唯一足够大以被归类为矮行星的天体，其直径约为 950 km；而灶神星、智神星和健神星的平均直径均小于 600 km$^{[3][4][5][6]}$。其余在矿物学上已被分类的天体尺寸一直延续到仅数米大小$^{[7]}$。小行星带中的物质分布极为稀疏，以至于多艘无人航天器已多次穿越该区域而未发生任何碰撞事故$^{[8]}$。尽管如此，大型小行星之间的碰撞仍然会发生，并能形成小行星族群，其成员具有相似的轨道特征和成分。带内各个小行星依据其光谱进行分类，大多数可归入三大基本类型：富碳的 C 型、富硅酸盐的 S 型以及富金属的 M 型。

小行星带是由原始太阳星云中的一群微行星形成的$^{[9]}$，这些微行星是原行星的较小前身。然而，在火星与木星之间，来自木星的引力扰动阻碍了这些天体进一步聚积成一颗行星$^{[9][10]}$，并赋予它们过量的动能，从而在碰撞时将微行星以及大多数初生原行星击碎。其结果是，在太阳系历史最初的一亿年内，小行星带原始质量的 99.9\% 被消耗殆尽$^{[11]}$。其中一些碎片最终进入太阳系内部区域，导致与类地行星的陨石撞击。只要小行星绕太阳的公转周期与木星形成轨道共振，其轨道就会持续受到显著扰动；在这些轨道距离处会形成所谓的柯克伍德空隙，因为小行星被“扫”入其他轨道上$^{[12]}$。
\subsection{观测史}
\begin{figure}[ht]
\centering
\includegraphics[width=6cm]{./figures/afb56775c2129f99.png}
\caption{1596 年，约翰内斯·开普勒（Johannes Kepler）基于自己对于行星轨道比例关系的直觉，认为在火星与木星的轨道之间存在一颗尚未被发现的行星$^{[13]}$。} \label{fig_belt_3}
\end{figure}
1596 年，约翰内斯·开普勒（Johannes Kepler）在其著作《宇宙秘典》（Mysterium Cosmographicum）中写道：“在火星与木星之间，我放置一颗行星”，明确表达了他认为在那里应当存在一颗行星的预测$^{[14]}$。在分析第谷·布拉赫（Tycho Brahe）的观测数据时，开普勒认为火星与木星轨道之间的间隙过大，不符合他关于行星轨道应当分布位置的模型$^{[15]}$。

在 1766 年他匿名发表的对夏尔·博内（Charles Bonnet）《自然沉思录》（Contemplation de la Nature）的译本附注中$^{[16]}$，维滕贝格的天文学家约翰·丹尼尔·提丢斯（Johann Daniel Titius）$^{[17][18]}$ 指出行星布局中似乎存在某种规律，这一规律后来被称为提丢斯–波得定律（Titius–Bode Law）。如果从 0 开始构造一个数列，然后依次取 3、6、12、24、48 等数（每次加倍），再给每个数字加上 4 并除以 10，就可以得到与当时已知行星轨道半径（以天文单位计）惊人接近的数值——前提是必须承认在火星（对应数列的 12）与木星（48）轨道之间存在一颗“缺失的行星”（对应数列的 24）。在附注中，提丢斯写道：“难道造物主会把那片空间空出来吗？绝不会。”$^{[17]}$

当威廉·赫歇尔（William Herschel）于 1781 年发现天王星时，其轨道与这一定律高度吻合，使得一些天文学家断定，在火星和木星轨道之间必然存在一颗行星$^{[19]}$。
\begin{figure}[ht]
\centering
\includegraphics[width=6cm]{./figures/7ecae5c9889c484d.png}
\caption{朱塞佩·皮亚齐（Giuseppe Piazzi），即小行星带中最大天体谷神星（Ceres）的发现者：谷神星曾被视为一颗行星，但随后被重新分类为小行星，并自 2006 年起被归类为矮行星。} \label{fig_belt_4}
\end{figure}
1801 年 1 月 1 日，西西里巴勒莫大学天文学主席朱塞佩·皮亚齐（Giuseppe Piazzi）发现了一个微小的移动天体，其轨道半径与该规律所预测的数值完全一致。他将其命名为“谷神星”（Ceres），取自罗马的丰收女神、亦为西西里的守护神。皮亚齐起初认为它可能是一颗彗星，但其缺乏彗发的特征表明它更像是一颗行星$^{[20]}$。由此，上述规律成功预测了当时八颗行星（即水星、金星、地球、火星、谷神星、木星、土星和天王星）的轨道半长轴。

与谷神星的发现同时，应弗朗茨·克萨韦尔·冯·扎克（Franz Xaver von Zach）的邀请，24 名天文学家组成了一个非正式团体，被称为“天上警察”（celestial police），其明确目的在于寻找更多行星；他们重点在火星与木星之间搜索，因提丢斯–波得定律预言该区应存在一颗行星$^{[21][22]}$。

大约 15 个月后，“天上警察”成员海因里希·奥伯斯（Heinrich Olbers）在同一区域发现了第二个天体——智神星（Pallas）。但与其他已知行星不同，谷神星与智神星在望远镜最高放大倍率下仍然呈现为光点，而非解析为圆面。除了快速移动外，它们与恒星几乎无差别$^{[23]}$。

因此，1802 年，威廉·赫歇尔（William Herschel）建议将它们划入一个独立类别，并以希腊语 *asteroeides*（意为“类星的”）为来源，为其命名为“asteroids”（小行星）$^{[24][25]}$。在完成对谷神星与智神星的一系列观测后，他得出结论$^{[26]}$：

无论是“行星”还是“彗星”这一称谓，都不能恰当地应用于这两个天体……它们与小恒星极为相似，以至于几乎无法将其区分开来。基于它们这种类星的外观，我取其名，称之为“小行星”；但我保留更改其名称的权利，若日后发现一个能更好表达其本质的术语。

到 1807 年，进一步观测又在该区域发现了两颗新天体：婚神星（Juno）与灶神星（Vesta）$^{[23]}$。然而，拿破仑战争中利林塔尔（Lilienthal）天文台的焚毁使这一早期发现阶段被迫终结$^{[27][23]}$。

尽管赫歇尔已提出“小行星”一词，但在接下来的几十年里，人们仍普遍将这些天体称为行星$^{[16]}$，并按发现顺序为其编号：1 谷神星、2 智神星、3 婚神星、4 灶神星。然而 1845 年天文学家卡尔·路德维希·亨克（Karl Ludwig Hencke）发现了第五个天体（5 义神星 Astraea），紧接着更多天体以加速的速度被发现。将它们都算作行星变得愈发不切实际。最终，它们被从行星名单中移除（最早由亚历山大·冯·洪堡在 1850 年代初提出），而赫歇尔所造的“小行星”（asteroids）一词逐渐成为通用称谓$^{[16]}$。

1846 年海王星（Neptune）的发现使提丢斯–波得定律在科学界迅速失去信誉，因为其轨道与该定律所预测的位置毫不相近。直至今日，这一定律仍无科学解释，天文学界普遍认为它只是巧合$^{[28]}$。
\begin{figure}[ht]
\centering
\includegraphics[width=6cm]{./figures/922f458ab47f5d40.png}
\caption{951 号小行星加斯普拉（Gaspra）——首个被航天器成像的小行星，其图像来自伽利略号（Galileo）在 1991 年飞掠期间拍摄；图像的颜色经过了夸张处理。} \label{fig_belt_5}
\end{figure}
“asteroid belt”（小行星带）这一表达在 19 世纪 50 年代初开始被使用，但确切由谁最先创造此术语已难以考证。其最早的英文用例似乎出现在亚历山大·冯·洪堡（Alexander von Humboldt）的《宇宙》（*Cosmos*）1850 年的英文译本（由 Elise Otté 翻译）中$^{[29]}$：“[…] 以及流星在 11 月 13 日与 8 月 11 日附近的规律出现，这些流星很可能属于一条与地球轨道相交、以行星速度运动的小行星带的一部分。”另一处早期的用例出现在罗伯特·詹姆斯·曼（Robert James Mann）的《天国知识指南》（*A Guide to the Knowledge of the Heavens*）中$^{[30]}$：“小行星的轨道位于一条宽广的空间带中，延伸于[…] 的极限之间。”美国天文学家本杰明·皮尔斯（Benjamin Peirce）似乎采用并推广了这一术语$^{[31]}$。

到 1868 年年中，已定位的小行星超过 100 颗，而 1891 年马克斯·沃尔夫（Max Wolf）将天体摄影技术引入小行星搜寻后，发现速度显著加快$^{[32]}$。到 1921 年已发现 1,000 颗小行星$^{[33]}$，1981 年达到 10,000 颗$^{[34]}$，并在 2000 年突破 100,000 颗$^{[35]}$。现代小行星巡天系统现已使用自动化方法，以不断增长的数量探测新的小行星。

2014 年 1 月 22 日，欧洲航天局（ESA）的科学家首次明确报告在小行星带最大天体谷神星（Ceres）上探测到水汽$^{[36]}$。该探测由赫歇尔空间天文台（Herschel Space Observatory）利用其远红外能力完成$^{[37]}$。这一发现出乎意料，因为通常认为会“喷出气流与羽状物”的是彗星而非小行星。正如其中一位科学家所言：“彗星与小行星之间的界线正变得越来越模糊。”$^{[37]}$
\subsection{起源}
\begin{figure}[ht]
\centering
\includegraphics[width=6cm]{./figures/4dd59d36fca97ed2.png}
\caption{小行星带示意图：展示了轨道倾角与距离太阳的关系，其中小行星带核心区域的天体以红色标示，其他小行星以蓝色标示。} \label{fig_belt_6}
\end{figure}
\subsubsection{形成}

1802 年，在发现智神星（Pallas）后不久，奥伯斯（Olbers）向赫歇尔（Herschel）和卡尔·高斯（Carl Gauss）提出：谷神星（Ceres）与智神星可能是曾经存在于火星—木星区域的一颗巨大行星的碎片，该行星在数百万年前因内部爆炸或彗星撞击而解体$^{[38]}$；而敖德山天文学家 K. N. Savchenko 则提出：谷神星、智神星、婚神星（Juno）和灶神星（Vesta）并非爆炸行星的碎片，而是“逃逸的卫星”$^{[39]}$。然而，摧毁一颗行星所需的能量极其巨大，加之小行星带的总质量仅约为地球月球质量的 4\%$^{[3]}$，均不支持这些假说。此外，小行星之间显著的化学差异也难以用同一母体行星来解释$^{[40]}$。

关于小行星带形成的现代假说与太阳系总体的行星形成过程相关，即与长期存在的星云假说相似：一团由星际尘埃与气体构成的云在引力作用下坍缩，形成一个旋转的物质盘，随后凝聚形成太阳与行星$^{[41]}$。在太阳系历史的最初几百万年里，微粒之间的黏性碰撞使其逐渐聚集成更大的团块；当这些团块的质量达到一定阈值后，它们便能够通过引力吸引其他物体，成为微行星（planetesimals）。这种引力聚积过程最终导致行星的形成$^{[42]}$。

但位于未来小行星带所在区域的微行星受到了木星引力的强烈扰动$^{[43]}$。轨道共振发生在小行星的公转周期与木星轨道周期构成整数比的区域，从而将小行星扰动至其他轨道；火星与木星之间存在大量此类轨道共振。当木星在形成后向内迁移时，这些共振会扫过整个小行星带，使其中天体的动力学激发程度上升，并提高其相互之间的相对速度$^{[44]}$。在某些区域，碰撞的平均速度过高，使得微行星更容易在碰撞中破碎而非聚积$^{[45]}$，从而阻止了一颗行星的形成。结果，它们继续绕太阳运行，并偶尔发生碰撞$^{[43]}$。

在太阳系早期历史中，小行星经历了一定程度的熔融，使得其内部物质发生按质量分异。一些母体天体甚至可能经历过阶段性的爆炸性火山活动，形成岩浆海。然而，由于这些天体体积较小，相比大型行星，其熔融期极为短暂，并普遍在约 45 亿年前、即形成后的最初数千万年内结束$^{[46]}$。2007 年 8 月，一项对据信源自灶神星的南极陨石中锆石晶体的研究发现，灶神星（进而小行星带）形成速度可能非常快——在太阳系形成后的 1000 万年内就已形成$^{[47]}$。
\subsubsection{演化}
\begin{figure}[ht]
\centering
\includegraphics[width=6cm]{./figures/c1ce78ad56022f73.png}
\caption{主小行星带的大型小行星——4 号灶神星（Vesta）} \label{fig_belt_7}
\end{figure}
小行星并不是原始太阳系的“原始样本”。它们自形成以来经历了显著的演化，包括内部加热（在最初几千万年内）、撞击造成的表面熔融、辐射造成的空间风化，以及微陨石的持续轰击$^{[48][49][50][51]}$。尽管一些科学家将小行星称为残余微行星（residual planetesimals）$^{[52]}$，但也有科学家认为它们已是独特的类群$^{[53]}$。

当前的小行星带被认为仅含有原始小行星带质量的一小部分。计算机模拟表明，原始小行星带的质量可能与地球相当$^{[54]}$。主要由于引力扰动，大部分物质在形成后的约 100 万年内被从小行星带中清除，仅剩不到原始质量的 0.1\%$^{[43]}$。自形成以来，小行星带的大小分布保持相对稳定；主带小行星的典型尺寸并未出现显著增大或减小$^{[55]}$。

与木星存在 4:1 轨道共振的位置（轨道半径 2.06 天文单位，AU）可视为小行星带的内边界。木星的引力扰动会使位于该处的天体被推入不稳定轨道。大多数形成在此间隙以内的天体在太阳系早期历史中被火星（其远日点为 1.67 AU）吸积，或被火星的引力扰动驱逐出该区域$^{[56]}$。洪伽利亚小行星族（Hungaria asteroids）位于 4:1 共振内侧更靠近太阳的位置，但由于其轨道倾角较高，得以避免被扰乱$^{[57]}$。

在小行星带形成初期，距离太阳约 2.7 AU 的位置形成了一个“雪线”（snow line），其温度低于水的冰点。形成于该半径以外的微行星能够积聚冰$^{[58][59]}$。2006 年，小行星带外侧雪线以外被发现存在一类彗星群，它们可能为地球海洋提供了水源。根据某些模型，地球形成早期由火山逸出所释放的水量不足以形成海洋，因此需要外部来源，如彗星轰击$^{[60]}$。

小行星带外部包含一些可能在过去几百年间被“植入”该区域的含冰天体。其中一个例子是类希尔达彗星（quasi-Hilda comet）362P/(457175) 2008 GO98，它被认为可能是一颗曾经的半人马型天体（centaur），通过与木星的近距离接触而被送入外小行星带$^{[61]}$。
\subsection{特性}
\begin{figure}[ht]
\centering
\includegraphics[width=6cm]{./figures/8cba3638ed75e30c.png}
\caption{主小行星带中小行星的大小分布$^{[62]}$} \label{fig_belt_8}
\end{figure}
\begin{figure}[ht]
\centering
\includegraphics[width=6cm]{./figures/b9ebd6f891ec6514.png}
\caption{} \label{fig_belt_9}
\end{figure}
与大众想象相反，小行星带大部分区域实际上是空旷的。小行星分布在极其巨大的体积中，因此如果不进行精确瞄准，想要接触到任意一颗小行星的概率非常低。然而，现今已知的小行星数量达数十万颗，总数可能达到数百万甚至更多，具体取决于最小尺寸的划定标准。已有超过 200 颗小行星被确认其尺寸大于 100 km$^{[63]}$，而红外波段的巡天研究显示，小行星带中直径 1 km 或以上的小行星数量介于 70 万至 170 万之间$^{[64]}$。

主小行星带中的小行星数量随其尺寸减小而稳定增加。虽然其尺寸分布总体上符合幂律关系，但在约 5 km 与 100 km 处存在“凸起”（bumps），即这些尺度的小行星比幂律预测的数量更多。大约 120 km 以上的多数小行星是原始天体，自吸积时代起便存续至今；而大部分较小的小行星则是原始天体破碎后的产物。主小行星带的原始天体族群数量可能是如今的 200 倍$^{[65][66]}$。

小行星之间的平均距离约为 965,600 km（600,000 英里）$^{[67][68]}$，尽管这一数值会随着不同小行星族群而变化，而较小且尚未被探测到的小行星之间的间距可能更小。小行星带的总质量估计为 (2.39 \times 10^{21}) kg，即月球质量的 3\%$^{[2]}$。四个最大天体——谷神星（Ceres）、灶神星（Vesta）、智神星（Pallas）与健神星（Hygiea）大约包含小行星带总质量的 62\%，其中仅谷神星一体就占 39\%$^{[69][5]}$。
\subsubsection{成分}
\begin{figure}[ht]
\centering
\includegraphics[width=8cm]{./figures/bd0444b740965155.png}
\caption{按距太阳距离划分的小行星光谱类型分布$^{[70]}$} \label{fig_belt_10}
\end{figure}
当今的小行星带主要由三大类小行星组成：C 型（碳质）小行星、S 型（硅酸盐）小行星，以及介于两者之间的混合族 X 型小行星。X 型混合族的小行星光谱无明显特征，但可根据反射率分为三类：M 型（金属质）、P 型（原始质），以及 E 型（顽辉石质）小行星。此外，还存在一些无法归入上述主要分类中的额外类型。小行星类型随其距太阳距离增加而呈现成分上的系统变化，依序为：S 型、C 型、P 型，及光谱无显著特征的 D 型小行星$^{[71]}$。
\begin{figure}[ht]
\centering
\includegraphics[width=6cm]{./figures/e1318718d025e78a.png}
\caption{阿连德陨石的碎片，这是一块富碳的球粒陨石，于1969年坠落在墨西哥。} \label{fig_belt_11}
\end{figure}
顾名思义，碳质小行星（C 型）富含碳。它们在小行星带的外部区域占主导地位$^{[72]}$，而在内带较为罕见$^{[71]}$。总体而言，C 型小行星占已观测到小行星总数的 75\% 以上。它们的色调比其他类型更偏红，反照率较低，其表面成分与碳质球粒陨石类似。从化学上看，它们的光谱与早期太阳系的原始成分一致，只是其中的氢、氦及挥发物已被移除$^{[73]}$。

S 型（硅酸盐丰富）小行星在距太阳 2.5 AU 以内的小行星带内侧区域更加常见$^{[72][74]}$。其表面光谱显示含有硅酸盐和少量金属，但缺乏显著的碳质化合物。这说明这些天体的物质已较多地偏离原始成分，可能经历了熔融与再形成过程。它们的反照率较高，占全部小行星的约 17\%$^{[73]}$。

M 型（金属质）小行星通常位于主带中部，构成了剩余小行星总体的重要部分$^{[73]}$。其光谱与铁–镍的光谱相似。有些被认为是经分异的母体天体的金属核，在碰撞中被破坏后所残存的部分。然而，一些硅酸盐化合物也可能产生类似的光谱。例如，大型 M 型小行星 22 号卡利俄珀（Kalliope）似乎并非主要由金属组成$^{[75]}$。在小行星带中，M 型小行星的数量分布在约 2.7 AU 的半长轴处达到峰值$^{[76]}$。目前尚不清楚所有 M 型小行星在成分上是否一致，或这一类别是否只是将若干无法清晰归入 C 型与 S 型的不同类型归为一组$^{[77]}$。

小行星带中的一个谜题是 V 型（Vestoid，灶神星族）或玄武岩质小行星的相对稀少$^{[78]}$。根据小行星形成理论，类似灶神星大小或更大的天体应当形成地壳与地幔，这些结构主要由玄武岩物质构成，因此应有超过一半的小行星由玄武岩或橄榄石组成。然而观测表明，预测的玄武岩物质中有 99\% 处于缺失状态$^{[79]}$。直到 2001 年为止，小行星带中发现的大多数玄武岩质天体都被认为起源于灶神星（故而得名 V 型），但小行星 1459 Magnya 的发现显示其化学成分与此前发现的其他玄武岩质小行星略有不同，暗示其可能具有不同的起源$^{[79]}$。这一假说在 2007 年得到进一步支持：在外小行星带发现了另外两颗玄武岩质小行星——7472 Kumakiri 与 (10537) 1991 RY16，它们的玄武岩成分与灶神星无法对应。截至目前，这两颗天体是唯一在外小行星带发现的 V 型小行星$^{[78]}$。
\begin{figure}[ht]
\centering
\includegraphics[width=6cm]{./figures/29d76c6abaccb71d.png}
\caption{哈勃望远镜观测到拥有多条尾巴的彗星状小行星 P/2013 P5。$^{[80]}$} \label{fig_belt_12}
\end{figure}
小行星带的温度随其距太阳的远近而变化。对于带内的尘埃粒子而言，典型温度范围从 2.2 AU 处的 200 K（−73 °C）下降到 3.2 AU 处的 165 K（−108 °C）$^{[81]}$。然而，由于自转的作用，小行星表面的温度可能出现显著变化，原因在于其不同侧面交替暴露于太阳辐射以及星际背景之中。
\subsubsection{主带彗星}
在外小行星带，一些原本毫不起眼的天体显示出彗星活动。由于它们的轨道无法用传统彗星的俘获机制来解释，因此人们认为外带的许多小行星可能含有冰，而这些冰会在受到小型撞击暴露后发生升华。主带彗星可能是地球海洋的重要来源，因为它们的氘氢比（D/H）远低于传统彗星，不足以支持后者为海水的主要来源$^{[82]}$。
\subsubsection{轨道}
\begin{figure}[ht]
\centering
\includegraphics[width=6cm]{./figures/df4a6e1bd847887f.png}
\caption{小行星带（显示偏心率），其中红色和蓝色部分代表小行星带（红色为‘核心’区域）。} \label{fig_belt_13}
\end{figure}
“小行星带内的大多数小行星轨道偏心率低于 0.4，轨道倾角低于 30°。小行星的轨道分布在偏心率约 0.07 和倾角低于 4° 处达到最大值。$^{[83]}$因此，尽管典型的小行星轨道相对圆且接近黄道平面，但也存在轨道高度偏心或远离黄道平面的情况。

有时，“主带”（main belt）一词专指最为紧凑、天体最为密集的‘核心’区域。该区域位于 2.06 和 3.27 天文单位处的强 4:1 与 2:1 柯克伍德缺口之间，其轨道偏心率通常小于约 0.33，轨道倾角低于约 20°。截至 2006 年，这一区域包含了太阳系内已发现并编号的小行星（小天体）总数的 93\%。$^{[84]}$JPL 小天体数据库现已列出超过 100 万颗已知的主带小行星。$^{[85]}$

\textbf{柯克伍德缺口}
\begin{figure}[ht]
\centering
\includegraphics[width=6cm]{./figures/b873154590100d05.png}
\caption{} \label{fig_belt_14}
\end{figure}
“小行星的半长轴用于描述其绕太阳轨道的尺寸，其数值决定了该小行星的轨道周期。1866 年，丹尼尔·柯克伍德宣布在这些天体绕日轨道距离的分布中发现了若干空隙。这些空隙位于它们的公转周期与木星公转周期成整数比的位置。柯克伍德提出，木星的引力扰动会将小行星从这些轨道上移除。$^{[86]}$

当小行星的平均轨道周期与木星的轨道周期成整数分数关系时，就会形成与这颗气体巨行星的平均运动共振，从而足以扰动小行星，使其进入新的轨道参数。原始的小行星之所以进入这些缺口，是因为木星轨道的迁移所致。$^{[87]}$ 此后，小行星主要因雅可夫斯基效应$^{[71]}$ 而迁移进入这些缺口轨道，也可能因引力扰动或碰撞进入。进入缺口后，小行星会在共振的作用下逐渐被推入一个更大或更小半长轴的随机新轨道。
\subsection{碰撞}
\begin{figure}[ht]
\centering
\includegraphics[width=6cm]{./figures/19205c038a5a95aa.png}
\caption{黄道光，其中一部分由行星际尘埃反射，而这些尘埃部分来源于小行星的碰撞。} \label{fig_belt_15}
\end{figure}
小行星带中极高的天体数量使其成为一个高度活跃的环境，其中小行星之间的碰撞在深时间尺度上频繁发生。平均半径为 10 公里的主带天体之间的撞击事件预计大约每一千万年发生一次$^{[88]}$。一次碰撞可能会将一颗小行星碎裂成许多更小的碎片（从而形成一个新的小行星族）$^{[89]}$。相反，以较低相对速度发生的碰撞也可能使两颗小行星结合。经过超过 40 亿年的此类演化过程，如今的小行星带成员与最初的族群已几乎毫无相似之处。

证据表明，直径介于 200 米到 10 公里之间的大多数主带小行星都是由碰撞形成的“碎石堆”。这些天体由大量形状不规则的碎块组成，主要依靠自引力聚合在一起，因此具有显著的内部孔隙度$^{[90]}$。除小行星本体外，小行星带中还包含粒径可达数百微米的尘埃带。这些细微物质部分来自小行星之间的撞击，以及微流星体撞击小行星产生的碎屑。由于彭廷–罗伯逊效应，太阳辐射压力会导致这些尘埃缓慢向太阳螺旋式内移$^{[91]}$。

这些细小的小行星尘埃与被抛出的彗星物质共同生成了黄道光。这种微弱的极光般辉光可在夜间自太阳方向沿黄道面观测到。产生可见黄道光的小行星尘埃粒子平均半径约 40 微米。主带黄道云粒子的典型寿命约为 70 万年。因此，为维持这些尘埃带，必须在小行星带内持续不断地产生新的尘埃粒子$^{[91]}$。曾经人们认为小行星碰撞是黄道光的主要来源之一。然而，Nesvorný 及其同事的计算机模拟显示，黄道光尘埃中有 85\% 来自木星家族彗星的碎裂，而非来自彗星与小行星带中小行星之间的碰撞。小行星带本身最多只贡献约 10\% 的尘埃$^{[92]}$。
\subsubsection{陨石}
部分碰撞碎片可能形成流星体并进入地球大气层$^{[93]}$。截至目前在地球上发现的 50,000 颗陨石中，有 99.8\% 被认为起源于小行星带$^{[94]}$。
\subsection{小行星族与群体}
\begin{figure}[ht]
\centering
\includegraphics[width=6cm]{./figures/075b2664084b42e1.png}
\caption{此图展示了编号主带小行星的轨道倾角（$(i_p)$）与偏心率（$(e_p)$）之间的关系，可以清晰地看到若干聚集区域，这些区域代表了不同的小行星族群。} \label{fig_belt_16}
\end{figure}
\begin{figure}[ht]
\centering
\includegraphics[width=6cm]{./figures/e8184c45bcc8665a.png}
\caption{内太阳系直到木星系统范围内的小行星概览} \label{fig_belt_17}
\end{figure}
\begin{figure}[ht]
\centering
\includegraphics[width=6cm]{./figures/a627147494f18831.png}
\caption{内太阳系天体的线性概览} \label{fig_belt_18}
\end{figure}
1918年，日本天文学家平山清次注意到一些小行星的轨道具有相似的参数，从而形成了所谓的“小行星族”或“小行星群”$^{[95]}$。

大约三分之一的主带小行星属于某个小行星族。它们共享相似的轨道要素，如半长轴、偏心率和轨道倾角，并且具有相似的光谱特征，这些都表明它们源自某个更大天体的解体。对主带小行星的这些轨道要素成对绘图时，可以看到若干明显的集中区，即小行星族的存在。目前已知大约有20–30个高度可能的小行星族，此外还发现了一些可信度较低的群体。当族群成员呈现相似的光谱特征时，小行星族即可得到确认$^{[96]}$。规模更小的聚集体通常被称为小行星群（groups）或小行星簇（clusters）。

“小行星带中最显著的一些族群（按半长轴从小到大排列）包括：Flora 族、Eunomia 族、Koronis 族、Eos 族以及 Themis 族$^{[76]}$。Flora 族是规模最大的族群之一，已知成员超过 800 个，可能形成于不足 10 亿年前的一次碰撞事件$^{[97]}$。在所有真正属于某一小行星族的成员中，体积最大的为 4 号灶神星（Vesta）。(与此相对，谷神星 Ceres 是 Gefion 族中的“闯入者”，并非真正成员。) Vesta 族被认为是一次在灶神星表面形成巨大撞击坑的事件所产生的。类似地，HED 陨石也可能正是这次撞击的直接产物$^{[98]}$。

在小行星带中发现了三条显著的尘埃带。它们的轨道倾角与 Eos、Koronis 和 Themis 小行星族相似，因此很可能与这些族群有关联$^{[99]}$。

在晚期重轰炸（Late Heavy Bombardment）之后，小行星带的演化很可能受到了大型半人马（centaur）天体与海王星外天体（TNOs）经过的影响。进入内太阳系的半人马与 TNOs 可以改变主带小行星的轨道——但仅限于它们的质量达到约 $1\times 10^{-9},M_{\odot}$（$2.0\times 10^{21}$ kg）时，单次掠过才足以产生影响；若为多次近距离掠过，则所需质量可低一个量级。尽管如此，半人马与 TNOs 不太可能显著扰乱主带中年轻的小行星族群，但它们确实可能影响部分年代久远的族群。当前主带中那些起源于半人马或 TNO 的小行星大多位于小行星带外部，其寿命通常不足 400 万年，并极可能以比典型主带小行星更高的偏心率，在 2.8–3.2 AU 之间运行$^{[100]}$。”
\subsubsection{边缘区域}
位于小行星带内侧边缘（范围约 1.78–2.0 AU，平均半长轴约 1.9 AU）的是匈牙利亚（Hungaria）小行星族。该族以其主成员 434 Hungaria 命名，包含至少 52 颗已命名小行星。匈牙利亚族与主带主体之间由 4:1 Kirkwood 缺口分隔，其轨道具有较高倾角。其中一些成员属于“火星穿越型小行星”，火星的引力扰动很可能是导致该族整体数量减少的重要因素$^{[57]}$。

主带内侧另一个高倾角族群为 Phocaea 族。该族主要由 S 型小行星组成，而邻近的 Hungaria 族中则包含部分 E 型小行星$^{[101]}$。Phocaea 族的轨道距离太阳约为 2.25–2.5 AU$^{[102]}$。

位于小行星带外缘的是 Cybele 族，其轨道范围约为 3.3–3.5 AU。Hilda 族的轨道位于 3.5–4.2 AU 之间，轨道较为圆形，并与木星保持稳定的 3:2 轨道共振。超过 4.2 AU 之后直到木星轨道之间的小行星数量非常稀少。在木星轨道附近则存在两大特洛伊小行星族群，若以直径大于 1 km 的天体计，其数量大致与整个小行星带相当$^{[103]}$。
\subsubsection{新形成的小行星族}
从天文学尺度来看，一些小行星族的形成相当年轻。例如，Karin 族似乎形成于约 570 万年前，其起源于一颗半径约 33 km 的母体小行星发生撞击碎裂$^{[104]}$。Veritas 族形成于约 830 万年前，相关证据包括从海洋沉积物中发现的行星际尘埃$^{[105]}$。

更为近期的例子是 Datura 族，推测约 53 万年前由一次主带小行星碰撞形成。其年龄估计基于成员拥有当前轨道的概率，而非直接物理证据；然而，这个族群可能是部分黄道光尘埃的来源$^{[106][107]}$。其它近期的小行星簇，例如 Iannini 簇（约 100–500 万年前），可能也为这些小行星尘埃提供了额外来源$^{[108]}$。
\subsection{探索}
\begin{figure}[ht]
\centering
\includegraphics[width=6cm]{./figures/0646c8ef89343ddf.png}
\caption{艺术家绘制的黎明号探测器及其与灶神星和谷神星的概念图} \label{fig_belt_19}
\end{figure}
第一艘穿越小行星带的航天器是先驱者10号（Pioneer 10），它于 1972 年 7 月 16 日进入该区域。当时有人担心小行星带中的碎片会对航天器构成危险，但之后多艘航天器均安全穿越而未发生事故。先驱者11号（Pioneer 11）、旅行者1号与2号（Voyager 1 and 2）以及尤利西斯号（Ulysses）都曾穿过小行星带，但未拍摄任何小行星影像。卡西尼号（Cassini）在 2000 年穿越小行星带时测量了等离子体和细微尘埃颗粒$^{[109]}$。在飞往木星的途中，朱诺号（Juno）穿越小行星带，但未采集科学数据$^{[110]}$。由于小行星带内部物质密度极低，探测器撞上小行星的概率估计低于十亿分之一$^{[111]}$。

迄今为止，大多数被成像的主带小行星都来自探测器在飞向其他目标时的短暂近距离飞掠机会。只有“黎明号”（Dawn）任务曾长期在轨研究主带小行星。伽利略号（Galileo）分别于 1991 年与 1993 年拍摄了 951 Gaspra 与 243 Ida；随后 NEAR 探测器于 1997 年成像 253 Mathilde，并于 2001 年 2 月登陆近地小行星 433 Eros。卡西尼号于 2000 年拍摄 2685 Masursky，星尘号（Stardust）于 2002 年拍摄 5535 Annefrank，“新视野号”（New Horizons）于 2006 年拍摄 132524 APL，“罗塞塔号”（Rosetta）则分别于 2008 年 9 月与 2010 年 7 月成像 2867 Šteins 与 21 Lutetia。“黎明号”在 2011 年 7 月至 2012 年 9 月环绕灶神星（Vesta）运行，并自 2015 年 3 月起一直环绕谷神星（Ceres）运行$^{[112]}$。

露西号（Lucy）探测器于 2023 年飞掠 152830 Dinkinesh，途中正前往木星特洛伊小行星群$^{[113]}$。欧洲航天局（ESA）的 JUICE 任务将两度穿越小行星带，并计划于 2029 年飞掠小行星 223 Rosa$^{[114]}$。Psyche 探测器则是 NASA 前往大型 M 型小行星 16 Psyche 的任务$^{[115]}$。”
\subsection{参见}
\begin{itemize}
\item Amor 小行星 —— 近地小行星的一类
\item Apollo 小行星 —— 近地小行星的一类
\item 小行星采矿 —— 对小行星原材料的开采与利用
\item Aten 小行星 —— 近地小行星的一类
\item 小行星带殖民 —— 关于人类殖民小行星带的拟议概念
\item 碎片盘（Debris disk）—— 围绕恒星运行的尘埃与碎片盘
\item 系外小行星（Exoasteroid）—— 太阳系外发现的小行星
\item 系外小行星带（Exoasteroid belt）
\item 特殊小行星列表（List of exceptional asteroids）
\item 小行星任务列表（List of missions to minor planets）
\item 费顿（Phaeton）—— 假想的太阳系原始行星
\end{itemize}
\subsection{参考资料}
\begin{enumerate}
\item 威廉姆斯，马特 (2015年8月23日)。“小行星带是什么？”。今天的宇宙。已检索1月30日，2016。
\item 皮特耶娃，E.V。(2018)。“来自行星和航天器运动的主要小行星带和柯伊伯带的质量”。太阳系研究。44(8-9):554-566.arXiv:1811.05191。Bibcode:2018年AstL...44 .. 554P。doi:10.1134/S1063773718090050。S2CID119404378。 
\item 克拉辛斯基，G.A.;皮特耶娃，E.V。；瓦西里耶夫，M。五、；Yagudina, E.I.(2002年7月)。“小行星带中的隐藏质量”。伊卡洛斯。158(1):98-105。Bibcode:2002年Icar .. 158...98K。doi:10.1006/icar.2002.6837。
\item 皮特耶娃，E.V。(2005)。“高精度行星星历表-EPM和某些天文常数的确定”(PDF)。太阳系研究。39(3):176-186。Bibcode:2005年索斯年。。39 .. 176便士。doi:10.1007/s11208-005-0033-2。S2CID120467483。存档自原来的(PDF)2014年7月3日  
\item 有关Ceres，Vesta，Pallas和Hygiea质量的最新估计，请参阅其各自文章的信息框中的参考资料。
\item 约曼，唐纳德·K。(2006年7月13日)。"JPL小体数据库浏览器"。NASA JPL。已存档从原来的2010年9月29日。已检索9月27日，2010。
\item “使用地球望远镜表征的已知最小小行星”。亚利桑那大学。2016年11月30日。已检索5月3日，2024。
\item Koberlein，Brian (2014年3月12日)。为什么小行星带不会威胁航天器。今天的宇宙。已检索1月30日，2016。
\item 小行星带是如何形成的？“那里有行星吗？”。CosmosUp.2016年1月17日存档自原来的2018年12月6日。已检索1月30日，2016。
\item 雷德，诺拉·泰勒 (2012年6月11日)。“小行星带: 事实与信息”。Space.com。已检索1月30日，2016。
\item Beatty，Kelly (2009年3月10日)。“雕刻小行星带”。天空和望远镜。已检索4月30日，2014。
\item 德尔格兰德，J。J.;索恩斯，S.V。(1943)。“柯克伍德在小行星轨道上的间隙”。加拿大皇家天文学会杂志。37: 187。Bibcode:1943JRASC..37..187D。
\item 坎宁安，克利福德J.(2022年9月)。“开普勒空白” 的起源和遗产”。天文历史与遗产杂志。25(3):439-456.Bibcode:2022JAHH...25..439C。doi:10.3724/SP.J.1440-2807.2022.03.02。S2CID 256570746。
\item 戴维斯，菲尔; 邓福德，比尔; 博克，摩尔。黎明: 在木星和火星之间，我放置了一颗行星(PDF)。喷气推进实验室。NASA。已存档(PDF)从2005年11月21日开始。
\item 罗素，克里斯托弗; 雷蒙德，卡罗尔，编辑。(2012)。黎明任务到小行星4 Vesta和1 Ceres。施普林格科学 + 商业媒体。第5页。ISBN 978-1-4614-4902-7。
\item 希尔顿，J。(2001)。“小行星是什么时候变成小行星的？”。美国海军天文台 (USNO)。存档自原来的2012年4月6日。已检索10月1日，2007。
\item 《黎明: 太阳系开端之旅》。空间物理中心: UCLA。2005。存档自原来的2012年5月24日。已检索11月3日，2007。
\item 霍斯金，迈克尔。“博德定律与谷神星的发现”。丘吉尔学院，剑桥。2008年5月10日从原件存档。已检索7月12日，2010。
\item 迈克尔·马丁·涅托 (2014)。Titius-bode行星距离定律及其历史和理论。爱尔斯维尔科学。第17页。ISBN 978-1-4831-5936-2。
\item “报警!小行星发现背后的故事 ”。天文学现在(2007年6月):60-61.
\item 温特伯恩，艾米丽 (2021年3月10日)。“发现小行星维斯塔: 天体警察的故事”。《夜空》杂志。英国广播公司。已检索10月18日，2022年。
\item 麦克威廉姆斯，布伦丹 (2001年7月31日)。“天体警察的徒劳搜索”。爱尔兰时报。已检索10月18日，2022年。
\item 工作人员 (2002年)。"天文意外"。NASA JPL。存档自原来的2012年2月6日。已检索4月20日，2007。
\item  哈珀，道格拉斯 (2010)。"小行星"。在线词源词典。词源在线。已检索4月15日，2011年。
\item 杰西卡·德福里斯特 (2000)。"希腊语和拉丁语根"。密歇根州立大学。已存档从原来的2007年8月12日。已检索7月25日，2007。
\item 坎宁安，克利福德 (1984)。“威廉·赫歇尔和前两个小行星”。小行星公报。11。安大略省舞厅天文台: 3。Bibcode:1984MPBu...11 .... 3C。
\item Elkins-Tanton, Linda T.(2011)。小行星、陨石和彗星(修订版)。纽约: 事实已存档。ISBN 978-0-8160-7696-3。OCLC 1054369860。[需要页面]
\item “大多数行星都在蒂修斯-博德定律的范围内，这是巧合吗？”。astronomy.com。已检索1月22日，2014。
\item  冯·洪堡，亚历山大 (1850)。宇宙: 宇宙的物理描述的草图。第一卷。纽约: 哈珀兄弟公司。第44页。ISBN 978-0-8018-5503-0。
\item 罗伯特·詹姆斯·曼 (1852)。天堂知识指南。贾罗尔德。第171页。和1853年，第216页
\item “有关小行星的形式，大小，质量和轨道的进一步研究”。爱丁堡新哲学杂志。5: 191。1857年1月至4月。: “[皮尔斯教授] 然后观察到土星环和小行星带之间的类比值得注意。”
\item 大卫·W·休斯.(2007)。“小行星发现的简要历史”。英国广播公司。存档自原来的2011年6月11日。已检索4月20日，2007。
\item 摩尔，帕特里克; 里斯，罗宾 (2011)。帕特里克·摩尔的天文学数据手册(第二版)。剑桥大学出版社。第156页。ISBN 978-0-521-89935-2。
\item Manley，Scott (2010年8月25日)。1980年至2010年的小行星发现。YouTube。存档自原来的2021年10月30日。已检索4月15日，2011年。
\item "MPC存档统计信息"。IAU小行星中心。已检索4月4日，2011年。
\item 迈克尔·库珀斯; 劳伦斯·奥罗克；Bockelée-Morvan, Dominique; Zakharov, Vladimir; Lee, Seungwon; von Allmen, Paul; Carry, Benoît; Teyssier, David; Marston, Anthony; Müller, Thomas; Crovisier, Jacques; Barucci, M.安东尼埃塔; 莫雷诺，拉斐尔 (2014)。“矮行星 (1) 谷神星上水蒸气的局部来源”。自然。505(7484):525-527。Bibcode:2014 Natur.505..525K。doi:10.1038/nature12918。ISSN0028-0836。PMID24451541。S2CID4448395。   
\item 哈林顿，J.D.(2014年1月22日)。“赫歇尔望远镜探测到矮行星上的水-发布14-021”。NASA。已检索1月22日，2014。
\item 坎宁安，克利福德J.(2017)。调查小行星的起源和Vesta的早期发现。小行星研究中的历史研究。施普林格国际出版。第1页。doi:10.1007/978-3-319-58118-7。ISBN 978-3-319-58117-0。
\item 布朗什顿，V。A.(1972)。“小行星的起源”。NASA.
\item 马塞蒂，M.；Mukai, K.(2005年12月1日)。“小行星带的起源”。NASA戈达德航天中心。已检索4月25日，2007。
\item 渡边，苏珊 (2001年7月20日)。“太阳星云的奥秘”。NASA.存档自原来的2012年1月17日。已检索4月2日，2007。
\item 钱伯斯，约翰E。(2004年7月)。“内太阳系的行星吸积量”。地球和行星科学快报。233(3-4):241-252.Bibcode:2004E & PSL.223..241C。doi:10.1016/j.epsl.2004.04.031。
\item J.-M小；莫比德利，A.；钱伯斯，J.(2001)。“小行星带的原始激发和清除”(PDF)。伊卡洛斯。153(2):338-347。Bibcode:2001年Icar .. 153 .. 338页。doi:10.1006/icar.2001.6702。已存档(PDF)从原来的2007年2月21日。已检索3月22日，2007。
\item 斯科特，E.R。D.(2006年3月13日至17日)。“球粒陨石和小行星对木星的年龄和形成机制以及星云寿命的限制”。第37届月球与行星科学年会论文集。德克萨斯州联盟城: 月球和行星协会。Bibcode:2006LPI....37.2367S。
\item 埃德加，R.；Artymowicz，P。(2004)。“通过快速迁移的行星泵送小行星盘”。皇家天文学会月刊。354(3):769-772。arXiv:astro-ph/0409017。Bibcode:2004MNRAS.354..769E。doi:10.1111/j.1365-2966.2004.08238.x。S2CID18355985。 
\item 泰勒，G。J.;凯尔，K.；麦考伊，T.；哈克，H; 斯科特，E。R。D.(1993)。“小行星分化 -- 火山碎屑火山作用到岩浆海洋”。气象学。28(1):34-52。Bibcode:1993年模拟 .. 28...34吨。doi:10.1111/j.1945-5100.1993.tb00247.x。
\item 凯利，凯伦 (2007)。“U of T研究人员发现了早期太阳系的线索”。多伦多大学。存档自原来的2012年2月29日。已检索7月12日，2010。
\item 克拉克，B。E.;哈普克，B.；皮特斯，C.；布里特，D。(2002)。“小行星空间风化与风化层演化”。小行星III。亚利桑那大学。第585页。Bibcode:2002aste.book .. 585C。doi:10.2307/j.ctv1v7zdn4.44。ISBN 978-0-8165-2281-1。
\item 加菲，迈克尔·J.(1996)。“陨石组合中金属的光谱和物理性质: 对小行星表面材料的影响”。伊卡洛斯。66(3):468-486。Bibcode:1986年Icar...66 .. 468克。doi:10.1016/0019-1035(86)90086-2。ISSN0019-1035。 
\item 凯尔，K。(2000)。“小行星的热蚀变: 来自陨石的证据”。行星与空间科学。48(10):887-903。Bibcode:2000P & SS。48。887k。doi:10.1016/S0032-0633(00)00054-4。已检索11月8日，2007。
\item 巴拉吉奥拉，R.A.;杜克，C。A.;洛夫勒，M.；麦克法登，L.A.;谢菲尔德，J。(2003)。“离子和微陨石对矿物表面的影响: 无空气太阳系体中大气物种的反射率变化和产生”。Egs-agu-eug联合组件: 7709。Bibcode:2003EAEJA.....7709B。
\item 查普曼，C。R、；威廉姆斯，J。G.;哈特曼，W。K。(1978)。“小行星”。天文学和天体物理学年度评论。16:33-75。Bibcode:1978ARA & A .. 16...33摄氏度。doi:10.1146/annurev.aa.16.090178.000341。
\item Kracher, A.(2005)。“小行星433 Eros和部分分化的小行星: 硫的体积消耗与表面消耗”(PDF)。艾姆斯实验室。已存档(PDF)从原来的2007年11月28日。已检索11月8日，2007。
\item Piccioni，罗伯特 (2012年11月19日)。“小行星撞击是否使地球宜居？”。Guidetothecosmos.com。已检索5月3日，2013。
\item Stiles，Lori (2005年9月15日)。“小行星造成了早期太阳系内部的灾难”。UANews。已检索10月18日，2018。
\item 阿尔夫文，H.; 阿伦尼乌斯，G。(1976)。《小身体》。太阳系的SP-345演化。NASA.已存档从原来的2007年5月13日。已检索4月12日，2007。
\item 斯普拉特，克里斯托弗E.(1990年4月)。“匈牙利小行星群”。加拿大皇家天文学会杂志。84:123-131.Bibcode:1990JRASC..84..123S。
\item 莱卡，M.；波多拉克，M.；萨瑟洛夫，D.；蒋，E。(2006)。“红外卷云-扩展红外发射的新组成部分”。天体物理学杂志。640(2):1115-1118。arXiv:astro-ph/0602217。Bibcode:2006ApJ。640.1115l。doi:10.1086/500287。S2CID18778001。 
\item Berardelli，Phil (2006年3月23日)。“主带彗星可能是地球水的来源”。每日空间。已检索10月27日，2007。
\item Lakdawalla，Emily (2006年4月28日)。“发现一种全新类型的彗星”。行星社会。存档自原来的2007年5月1日。已检索4月20日，2007。
\item de la Fuente Marcos，卡洛斯; de la Fuente Marcos，劳尔 (2022年10月1日)。“最近到达主要小行星带的人”。天体力学与动力天文学。134(5): 38。arXiv:2207.07013。Bibcode:2022年CeMDA.134. .. 38D。doi:10.1007/s10569-022-10094-4。
\item 戴维斯，D。R、；杜达，D.D.;马尔扎里，F.；Campo Bagatin, A.;吉尔-赫顿，R。(2002年3月)。“小体种群的碰撞演化”(PDF)。在Bottke Jr.，W.F.;切利诺，A.；Paolicchi, P.;宾泽尔，R。P、 (编辑)。小行星III。图森: 亚利桑那大学出版社。pp. 545-558。Bibcode:2002年aste.book .. 545D。doi:10.2307/j.ctv1v7zdn4.41。ISBN  978-0-8165-2281-1。已检索2月11日，2023。
\item 约曼斯，唐纳德·K。(2007年4月26日)。"JPL小型数据库搜索引擎"。NASA JPL。已检索4月26日，2007。 -在直径> 100的主带区域中搜索小行星。
\item 特德斯科，E。F.;沙漠，F.-X。(2002)。“红外空间天文台深小行星搜索”。天文杂志。123(4):2070-2082。Bibcode:2002AJ....123.2070T。doi:10.1086/339482。
\item Bottkejr, W.;杜尔达，D.；内斯沃尼，D.；Jedicke, R.;莫比德利，A.；Vokrouhlicky, D.;Levison，H. (2005年5月)。“主要小行星带的化石大小分布”。伊卡洛斯。175(1):111-140.Bibcode:2005年Icar。。175。。111B。doi:10.1016/j.icarus.2004.10.026。
\item 大卫·P·奥布莱恩；Sykes, Mark V.(2011年12月)。“小行星带的起源和演化-对Vesta和Ceres的影响”。空间科学评论。163(1-4):41-61.Bibcode:2011SSRv。。163...41O。doi:10.1007/s11214-011-9808-6。ISSN0038-6308。S2CID121856071。  
\item “小行星带包含太阳系残余物”。EarthSky | 更新你的宇宙和世界。2021年11月3日。已检索1月20日，2023。
\item 威廉姆斯，马特 (2016年8月10日)。“小行星带离地球有多远？”。今天的宇宙。已检索1月20日，2023。
\item “深度 | Ceres”。NASA太阳系探索。2017年11月9日。已检索2月11日，2023。
\item 格拉迪，J.；特德斯科，E。(1982年6月)。“小行星带的组成结构”。科学。216(4553):1405-1407。Bibcode:1982Sci...216.1405克。doi:10.1126/science.216.4553.1405。PMID17798362。S2CID32447726。  
\item DeMeo, F.E.;亚历山大，C。M。奥德。；沃尔什，K.J.;查普曼，C。R、；宾泽尔，R。P。(2015)。“小行星带的组成结构”。在米歇尔，帕特里克; 德梅奥，弗朗西斯卡E.；威廉·F·博克 (编辑)。小行星IV。图森: 亚利桑那大学出版社。pp. 13-41.arXiv:1506.04805。Bibcode:2015aste.book...13D。doi:10.2458/azu_uapress_9780816532131-ch002。ISBN 978-0-816-53213-1。S2CID 4648806。
\item 维格特，P.；巴拉姆，D.；莫斯，A.；Veillet, C.;康纳斯，M.；谢尔顿，我(2007)。“主带小行星大小分布中颜色依赖性的证据”(PDF)。天文杂志。133(4):1609-1614.arXiv:astro-ph/0611310。Bibcode:2007AJ....133.1609W。doi:10.1086/512128。S2CID54937918。已存档(PDF)从2011年7月6日起。已检索9月6日，2008。  
\item 布莱尔，爱德华·C，编者 (2002年)。小行星: 概述、摘要和参考书目。新星科学出版社。第2页。ISBN 978-1-59033-482-9。
\item 克拉克，B。E.(1996)。“小行星带地质学的新新闻和相互竞争的观点”。月球和行星科学。27:225-226。Bibcode:1996LPI....27..225C。
\item 玛戈特，J.L.;布朗，M.E.(2003)。“主带中的低密度M型小行星”(PDF)。科学。300(5627):1939年-1942.Bibcode:2003Sci...300.1939米。doi:10.1126/science.1085844。PMID12817147。S2CID5479442。存档自原来的(PDF)2020年2月26日。   
\item  朗，肯尼斯·R。(2003)。“小行星和陨石”。NASA的宇宙。存档自原来的2012年3月24日。已检索4月2日，2007。
 穆勒，M.；哈里斯，A.W.;德尔博，M.(2005)。“21 Lutetia和其他M型: 它们的大小、反照率和来自新IRTF测量的热特性”。美国天文学会公报。37。MIRSI团队: 627。Bibcode:2005DPS....37.0702M。
 Duffard, R.D.;罗伊格，F。(2008年7月14日至18日)。“主带中有两个新的玄武岩小行星？”。小行星，彗星，流星2008。马里兰州巴尔的摩。arXiv:0704.0230。Bibcode:2008LPICo1405.8154D。
 比，Ker (2007)。“奇怪的小行星挡板科学家”。space.com。已检索10月14日，2007。
 “彗星什么时候不是彗星？”。ESA/哈勃新闻稿。已检索11月12日，2013。
 低，F。J.;贝因特玛，D。A.;戈蒂埃，T.N.;吉列特，F。C.;贝希曼，C。A.;纽格鲍尔，G.；杨，E.；奥曼，H.；Boggess, N.；艾默生，J。P.;哈宾，h.j.；豪瑟，M。G.;霍克，J。R、；罗文-罗宾逊，M.；Soifer, B.T.;沃克，R。G.;韦塞利厄斯，P。R。(1984)。“红外卷云-扩展红外发射的新组成部分”。天体物理学杂志快报。278:L19-L22。Bibcode:1984ApJ...278L..19L。doi:10.1086/184213。
 “采访David Jewitt”。YouTube.com。2007年1月5日存档自原来的2021年10月30日。已检索5月21日，2011年。
 威廉姆斯，加雷斯 (2010年9月25日)。“小行星的分布”。小行星中心。已检索10月27日，2010。
 该值是通过使用来自该地区的120，437个编号的小行星的数据对该地区的所有天体进行简单计数而获得的。小行星中心轨道数据库,日期为2006年2月8日。
 “JPL小体数据库搜索引擎: 轨道类 (MBA)”。JPL太阳能系统动力学。已检索2月26日，2018。
 Fernie, J.唐纳德 (1999)。《美国开普勒》。美国科学家。87(5): 398。doi:10.1511/1999.5.398(2025年7月1日不活跃)。存档自原来的2017年6月21日。已检索2月4日，2007。
 Liou，jer-chyi; Malhotra，Renu (1997)。“小行星带外带的枯竭”。科学。275(5298):375-377。Bibcode:1997Sci...275 .. 375L。doi:10.1126/science.275.5298.375。高密度脂蛋白:2060/19970022113。PMID8994031。S2CID33032137。  
 Backman, D.E.(1998年3月6日)。“一般黄道云密度的波动”。Backman报告。NASA艾姆斯研究中心。存档自原来的2012年3月3日。已检索4月4日，2007。
 内斯沃尼，大卫; 小博特克，威廉·F。；卢克·唐斯; 哈罗德·F·莱维森。(2002年6月)。“最近一颗小行星在主带地区的解体”(PDF)。自然。417(6890):720-722。Bibcode:2002 Natur.417..720N。doi:10.1038/nature00789。PMID12066178。S2CID4367081。已存档(PDF)从2004年8月7日开始。   
 凯文·J·沃尔什.(2018年9月)。“瓦砾堆小行星”。天文学和天体物理学年度评论。56:593-624。arXiv:1810.01815。Bibcode:2018ARA & A .. 56 .. 593W。doi:10.1146/annurev-astro-081817-052013。S2CID119530506。 
 到达，威廉·T。(1992)。“黄道带辐射。III-小行星带附近的尘埃“。天体物理学杂志。392(1):289-299。Bibcode:1992年4月。392 .. 289R。doi:10.1086/171428。
 内斯沃尼，大卫; 詹尼斯肯斯，彼得; 利维森，哈罗德·F。；威廉·F·博克；Vokrouhlick ý，大卫; Gounelle，Matthieu (2010)。“黄道带云和碳质微陨石的彗星起源。对热碎片盘的影响 “。天体物理学杂志。713(2):816-836。arXiv:0909.4322。Bibcode:2010年4月。713 .. 816N。doi:10.1088/0004-637X/713/2/816。S2CID18865066。 
 金斯利，丹尼 (2003年5月1日)。“神秘的陨石尘埃不匹配解决”。ABC科学。已检索4月4日，2007。
 “流星和陨石”(PDF)。NASA.存档自原来的(PDF)2006年10月15日。已检索1月12日，2012。
 大卫·W·休斯.(2007)。“在太空中寻找小行星”。英国广播公司。存档自原来的2012年3月10日。已检索4月20日，2007。
 Lemaitre，安妮 (2004年8月31日-9月4日)。“来自大型目录的小行星家族分类”。行星系统种群的动力学。贝尔格莱德，塞尔维亚和黑山: 剑桥大学出版社。pp. 135-144.Bibcode:2005dpps.conf .. 135L。doi:10.1017/S1743921304008592。
 Martel，琳达M.V。(2004年3月9日)。“一个大小行星破裂的微小痕迹”。行星科学研究发现。已存档从原来的2007年4月1日。已检索4月2日，2007。
 德雷克，迈克尔·J.(2001)。“eucrite/Vesta故事”。气象学与行星科学。36(4):501-513。Bibcode:2001M & PS...36 .. 501D。doi:10.1111/j.1945-5100.2001.tb01892.x。
 爱，S.G.;布朗利，D。E.(1992)。“IRAS尘埃带对行星际尘埃复合体的贡献-在60和100微米处看到的证据”。天文杂志。104(6):2236-2242.Bibcode:1992AJ....104.2236L。doi:10.1086/116399。
 Galiazzo, M.A.;维格特，P.；Aljbaae, S.(2016)。“半人马和TNOs对主带及其家族的影响”。天体物理学和空间科学。361(12):361-371。arXiv:1611.05731。Bibcode:2016Ap & SS.361..371G。doi:10.1007/s10509-016-2957-z。S2CID118898917。 
 卡瓦诺，J.M、；拉扎罗，D.；Mothé-Diniz, T.;Angeli, C.A.;Florczak, M.(2001)。“匈牙利和Phocaea动力组的光谱调查”。伊卡洛斯。149(1):173-189.Bibcode:2001年Icar。。149。。173C。doi:10.1006/icar.2000.6512。
 诺瓦科维奇，博扬; 齐尔沃利斯，乔治; 格兰维克，米凯尔; 托多维奇，安娜 (2017年6月)。“Phocaea地区的一个黑暗小行星家族”。天文杂志。153(6): 266。arXiv:1704.06088。Bibcode:2017AJ....153..266N。doi:10.3847/1538-3881/aa6ea8。ISSN0004-6256。S2CID96428710。  
 Dymock，Roger (2010)。小行星和矮行星以及如何观察它们。天文学家的观察指南。斯普林格。第24页。ISBN 978-1-4419-6438-0。已检索4月4日，2011年。
 内斯沃尼，大卫; 恩克，布莱恩L.；威廉·F·博克；丹尼尔·D·杜尔达；阿斯豪格，埃里克; 理查森，德里克·C。(2006年8月)。“小行星撞击形成卡林星团”。伊卡洛斯。183(2):296-311。Bibcode:2006年Icar。。183。。296N。doi:10.1016/j.icarus.2006.03.008。
 麦基，玛吉 (2006年1月18日)。“沙尘暴的Eon可追溯到小行星粉碎”。新科学家空间。已检索4月15日，2007。
 Nesvorný, D.;Vokrouhlický, D.;Bottke, W.F.(2006)。“45万年前一颗主带小行星的解体”(PDF)。科学。312(5779): 1490。Bibcode:2006Sci。312.1490n。doi:10.1126/science.1126175。PMID16763141。S2CID38364772。已存档(PDF)从2008年5月9日开始。   
 Vokrouhlický, D.;杜雷奇，J.；Michałowski, T.;克鲁格，于。N.;Gaftonyuk, N.M、；Kryszczyńska, A.;可乐，F.；Lecacheux, J.;莫洛托夫，I.；Slyusarev, I.;波利斯卡，M.；Nesvorný, D.;贝索尔，E。(2009)。“曼陀罗家族: 2009年更新”。天文学与天体物理学。507(1):495-504。Bibcode:2009A & A。507 .. 495伏。doi:10.1051/0004-6361/200912696。
 Nesvorný, D.;Bottke, W.F.;利维森，h.f.；Dones, L.(2003)。“太阳系尘埃带的最新起源”(PDF)。天体物理学杂志。591(1):486-497.Bibcode:2003ApJ...591 .. 486N。doi:10.1086/374807。S2CID1747264。存档自原来的(PDF)于2020年4月12日。  
 Schippers, P.;Meyer-Vernet, N.;Lecacheux, A.;Belheouane, S.;蒙库克，M.；Kurth, W.美国；曼恩，I.；米切尔，D。G.;安德烈，N.(2015年6月)。“使用卡西尼波测量在1到5 AU之间检测纳米灰尘”。天体物理学杂志。806(1): 77。arXiv:1504.02345。Bibcode:2015年apj... 806... 77s。doi:10.1088/0004-637X/806/1/77。S2CID118554353。77. 
 Greicius，托尼 (2015年7月31日)。“NASA的朱诺号为地球飞越提供了类似星舰的视野”。nasa.gov。NASA。已检索9月4日，2015。
 斯特恩，艾伦 (2006年6月2日)。“新视野号穿越小行星带”。每日空间。已检索4月14日，2007。
 巴鲁奇，M.A.;富基尼诺尼，M.；罗西，A.(2007)。“罗塞塔小行星目标: 2867斯坦斯和21卢特蒂亚”。空间科学评论。128(1-4):67-78。Bibcode:2007SSRv .. 128...67B。doi:10.1007/s11214-006-9029-6。S2CID123088075。 
 NASA的露西团队宣布新的小行星目标。NASA.2023年1月25日。已检索2月14日，2023。
 Avdellidou, C.;帕乔拉，M.；卢切蒂，A.；阿戈斯蒂尼，L.；德尔博，M.；Mazzotta Epifani, E.;Bourdelle De Micas, J.;德沃格勒，M.；Fornasier, S.;范·贝尔，G.；布鲁特，N.；Dotto, E.;Ieva, S.;克雷蒙人，G.；帕伦博，P。(2021)。“主带小行星 (223) 罗莎的特征”。天文学与天体物理学。656: l18。Bibcode:2021A & A...656L..18A。doi:10.1051/0004-6361/202142600。高密度脂蛋白:11577/3460828。S2CID244753425。 
 “NASA继续执行Psyche小行星任务”。NASA.2022年10月28日。已检索2月14日，2023。
 \end{enumerate}